\begin{center}
\hypertarget{native-059-article041}{}\pdfbookmark[2]{41.}{article-041}41.
\end{center}
Nihil iam superest, nisi ut etiam problema inversum, puta determinationem temporis ex anomalia vera, ad algorithmum expeditiorem reducamus: ad hunc finem tabulae nostrae columnam novam pro $T$ adiecimus. Computabitur itaque primo $T$ ex $v$ per formulam $T=\frac{1-e}{1+e}\operatorname{tang}\tfrac{1}{2}v^2$; dein ex tabula nostra argumento $T$ desumetur $A$ et $\log B$, sive (quod exactius, imo etiam commodius est) $C$ et $\log B$, atque hinc $A$ per formulam $A=\frac{(1+C)T}{1+\frac{4}{5}T}$; tandem ex $A$ et $B$ eruetur $t$ per formulam 1) art. 37. Quodsi hic quoque tabulam \textsc{Barker}ianam in usum vocare placet, quod tamen in hoc problemate inverso calculum minus sublevat, non opus est ad $A$ respicere, sed statim habetur
\[
\operatorname{tang}\tfrac{1}{2}w=\operatorname{tang}\tfrac{1}{2}v\,\mathord{.}\,\BareRoot{\frac{1+C}{\gamma\gamma(1+\frac{4}{5}T)}},
\]
atque hinc tempus $t$, multiplicando motum medium anomaliae verae $w$ in tabula \textsc{Barker}iana respondentem per $\frac{B}{\alpha}\cdot$
